% ============================================================
%  Buzzer Sub-Circuit  –  Revised layout
%  RPi GPIO → R1 → MMBT3904 (NPN) → Buzzer (vertical) → +5V
%  Flyback diode 1N4148W to the left of the buzzer
% ============================================================

\documentclass[border=12pt]{standalone}
\usepackage{circuitikz}
\usepackage{xcolor}

% ── Layout constants (easy to tweak) ──────────────────────
%  y=0   : R1 / NPN base line
%  x=0   : left edge (GPIO label)
%  x=6   : NPN centre
%  x=9   : buzzer / diode column

\begin{document}
	\begin{circuitikz}[american, scale=1.15, transform shape,
		/tikz/circuitikz/bipoles/length=1.4cm]
		% bipoles/length scales component body size uniformly
		
		% ══════════════════════════════════════════════════════
		%  TRANSISTOR  (centre of universe)
		% ══════════════════════════════════════════════════════
		% npn node centred at (6, -1) so that the Base pin
		% sits at y=0 (base is ~1 unit above centre for npn)
		\node[npn] (T1) at (6, -1) {};
		
		% ══════════════════════════════════════════════════════
		%  INPUT PATH:  GPIO  ──  R1  ──  Base
		% ══════════════════════════════════════════════════════
		% Open-circle pin (o-) at left, plain wire (-) at right
		\draw (0, 0)
		node[left]{\textbf{GPIO (PWM)}}
		to[R, o-, l=R1, a=1\,k$\Omega$]   % resistor on the wire
		(T1.B);                             % ends at Base pin
		
		% ══════════════════════════════════════════════════════
		%  EMITTER  ──  GND
		% ══════════════════════════════════════════════════════
		\draw (T1.E)
		to[short]
		(T1.E |- 0,-3)      % straight down to y=-3
		node[ground]{};     % GND symbol
		
		% ══════════════════════════════════════════════════════
		%  COLLECTOR  ──  bottom of buzzer
		% ══════════════════════════════════════════════════════
		% Named coordinate where collector wire meets buzzer column
		\coordinate (CBOT) at (9, -1);   % buzzer bottom terminal
		\coordinate (CTOP) at (9,  2);   % buzzer top  terminal  (+5V)
		
		% Horizontal wire: Collector → buzzer-bottom
		\draw (T1.C) -- (CBOT);
		
		% ══════════════════════════════════════════════════════
		%  BUZZER  (vertical, bottom→top = Collector→+5V)
		% ══════════════════════════════════════════════════════
		% CircuiTikZ 'buzzer' draws horizontally by default;
		% 'rotate=-90' inside component options re-orients it
		% so current flows bottom-to-top.
		% '*-*' puts filled junction dots at both terminals.
		\draw (CBOT)
		to[buzzer, *-*, rotate=-90,           % vertical orientation
		l_={\small CEM-1203(42)},           % label to right (l_ = right side)
		a_={\small Piezo Buzzer}]           % annotation to right
		(CTOP);
		
		% ══════════════════════════════════════════════════════
		%  +5V  RAIL
		% ══════════════════════════════════════════════════════
		\draw (CTOP)
		to[short, -*]              % short wire, dot at end
		++(0, 0.3)                 % tiny upward stub
		node[vcc]{\textbf{+5V}};  % VCC / supply symbol
		
		% ══════════════════════════════════════════════════════
		%  FLYBACK DIODE  (left of buzzer, anode at bottom)
		% ══════════════════════════════════════════════════════
		% Placed at x=7.2, spanning same y range as buzzer.
		% Anode (positive) at bottom → Collector side (CBOT y)
		% Cathode at top             → +5V side       (CTOP y)
		% Drawing path goes bottom→top, so normal 'D' orientation
		% has anode at start = bottom. ✓
		\coordinate (DBOT) at (7.2, -1);   % anode  (= CBOT y)
		\coordinate (DTOP) at (7.2,  2);   % cathode (= CTOP y)
		
		% Horizontal connecting wires from buzzer terminals to diode
		\draw (CBOT) -- (DBOT);   % bottom junction
		\draw (CTOP) -- (DTOP);   % top    junction
		
		% Diode drawn bottom→top; labels forced to the left
		\draw (DBOT)
		to[D, *-*,                           % dots at both ends
		l_={\small 1N4148W},              % label  left
		a_={\small Flyback}]              % annotation left
		(DTOP);
		% Note: l_ and a_ with the underscore push labels to the
		% "negative" side of the path — which is the left when
		% the path runs upward.
		
		% ══════════════════════════════════════════════════════
		%  TERMINAL LABELS  (minimal)
		% ══════════════════════════════════════════════════════
		\node[below right, font=\tiny] at (T1.B) {B};
		\node[above right, font=\tiny] at (T1.C) {C};
		\node[below right, font=\tiny] at (T1.E) {E};
		\node[below, font=\small]      at (T1.center) {MMBT3904};
		
	\end{circuitikz}
\end{document}

